\documentclass[11pt]{article}
\usepackage[utf8]{inputenc} 
\usepackage{a4wide} 
\usepackage[russian]{babel} 
\usepackage{graphicx} 
\usepackage{amsmath} 
\usepackage{amsfonts} 
\usepackage{amssymb} 
\usepackage{amsthm}
\usepackage{float}

\newtheorem{note}{Замечание}[section]
\newtheorem{definition}{Определение}[section] 
\newtheorem{theorem}{Теорема}[section] 
\newtheorem{property}{Свойство}[section] 
\newtheorem{corollary}{Следствие}[theorem] 
\newtheorem{lemma}[theorem]{Лемма} 
\newtheorem*{statement}{Утверждение} 
\renewenvironment{proof}{
\noindent\textit{Доказательство: }} {\qed}



\begin{document}
	%\thispagestyle{empty}
\begin{center}
    \ \vspace{-3cm}
    
    \includegraphics[width=0.5
    \textwidth]{msu}\\
    {\scshape Московский государственный университет имени М.~В.~Ломоносова}\\
    Факультет вычислительной математики и кибернетики\\
    Кафедра системного анализа
    
    \vfill
    
    {\LARGE Отчёт по практикуму}
    
    \vspace{1cm}
    
    {\Huge\bfseries "<Нелинейная задача оптимального управления">\\ (задача управления тележкой)} 
\end{center}
\vspace{1cm}
\begin{flushright}
    \large \textit{Студент 315 группы}\\
    В.\,В.~Кожемяк
    
    \vspace{5mm}
    
    \textit{Руководители практикума}\\
    к.ф.-м.н., доцент П.\,А.~Точилин 
\end{flushright}

\vfill
\begin{center}
    Москва, 2018 
\end{center}
\pagebreak

%Contents
\tableofcontents

\pagebreak


\section{Постановка задачи}
	Движение материальной точки на прямой описывается обыкновенным дифференциальным уравнением:
	$$
		\ddot{x} = u_1 - \dot{x}(1 + u_2), \; t \in [0, T],
	$$
	где $ x \in \mathbb{R}, \; u = (u_1, u_2)^T \in \mathbb{R}^2. $ На возможные значения управляющих параметров $ u_1, u_2 $ наложены следующие ограничения:
	\begin{enumerate}
		\item либо $ u_1 \geqslant 0, \; u_2 \in [0, k], \; k > 0, $
		\item либо $ u_1 \in \mathbb{R}, \; u_2 \in [0, k], \; k > 0. $
	\end{enumerate}
	Задан начальный момент времени $ t_0 = 0. $ Начальная позиция $ x(0) $ может быть выбрана произвольной, но так, чтобы выполнялось условие $ x(0) \in [x_{min}, x_{max}]. $ В начальный момент времени материальная точка неподвижна: $ \dot{x}(0) = 0. $ Необходимо за счёт выбора программного управления $ u $ и начальной позиции $ x(0) $ перевести материальную точку из начальной позиции в такую позицию в момент времени $ T, $ в которой $ x(T) = L > x_{max}, \; \vert \dot{x}(T) \vert \leqslant \varepsilon. $ На множестве всех реализаций программных управлений, переводящих материальную точку в указанное множество, необходимо решить задачу оптимизации:
	$$
		J = \int \limits_0^T u_1^2(t) dt \rightarrow \min \limits_{ u(\cdot) }
	$$
	\begin{enumerate}
		\item Необходимо написать в среде \texttt{Matlab} программу с пользовательским интерфейсом, которая по заданным параметрам $ T, k, x_{min}, x_{max}, L, \varepsilon $ определяет, разрешима ли задача оптимального управления (при одном из указанных двух ограничений на управления). Если задача разрешима, то программа должна построить графики компонент оптимального управления, компонент оптимальной траектории, сопряжённых переменных. Кроме того, программа должна определить количество переключений найденного оптимального управления, а также моменты переключений.
		\item В соответствующем задания отчёте необходимо привести все теоретические выкладки, сделанные в ходе решения задачи оптимального управления, привести примеры построенных оптимальных управлений и траекторий (с иллюстрациями). Все вспомогательные утверждения (за исключением принципа максимума Понтрягина), указанные в отчёте, должны быть доказаны. В отчёте должны быть приведены примеры оптимальных траекторий для всех возможных качественно различных ``режимов''.
	\end{enumerate}
	\textbf{Замечание.} Алгоритм построения оптимальной траектории не должен содержать перебра по параметрам, значения которых не ограничены, а также по более чем двум скалярным параматрам.
	
\section{Анализ решения}
	Обозначим $ x = x_1, \; \dot{x} = x_2 $ и запишем обыкновенное дифференциальное уравнение второго порядка:
	$$
		\ddot{x} = u_1 - \dot{x}(1 + u_2), \; t \in [0, T],
	$$
	в виде системы обыкновенных дифференциальных уравнений первого порядка
	\begin{equation}
		\label{odeSystem}
		\begin{cases}
   			\dot{x_1} = x_2, \\ 
   			\dot{x_2} = u_1 - x_2(1 + u_2),
 		\end{cases}
	\end{equation}
	с краевыми условиями
	\begin{equation*}
		\begin{cases}
   			x_1(0) = \alpha, \\
   			x_2(0) = 0, \\
   			x_1(T) = L, \\
   			\vert x_2(T) \vert \leqslant \varepsilon,
 		\end{cases}
	\end{equation*}
	где $ x_{min}, x_{max} $ --- некоторые заданные числа и $ \alpha \in [x_{min}, x_{max}]. $ \\
	Рассмотрим уравнение 
	$$
		\dot{x_2} = u_1 - x_2(1 + u_2). 
	$$
	Решение данного уравнения даёт формула Коши:
	$$
		x_2(t) = \chi(t, t_0)x_2(0) + \int \limits_{t_0}^t \chi(t, \tau)f(\tau) d\tau,
	$$
	где $ \chi(t, \tau) = \exp \left(- \int \limits_\tau^t (1 + u_2(s)) ds \right). $ \\
	Подставляя $ f(\tau) = u_1, x_2(0) = 0 $ в $ x_2(t), $ полчаем
	\begin{equation}
	\label{x2(t)}
		x_2(t) = \int \limits_{t_0}^t u_1(\tau) \cdot \exp \left(- \int \limits_\tau^t (1 + u_2(s)) ds \right) d\tau.
	\end{equation}
	Далее запишем функцию Гамильтона-Понтрягина
	$$
		H = \psi_0 u_1^2 + \psi_1 x_2 + \psi_2 (u_1 - x_2(1 + u_2))
	$$
	и сопряжённую систему
	\begin{equation}
		\label{conjSystem}
		\begin{cases}
			\dot{\psi_1} = 0, \\
			\dot{\psi_2} = - \psi_1 + \psi_2(1 + u_2).
   		\end{cases}
	\end{equation}
	Так как $ \dot{\psi_1} = 0, $ следовательно $ \psi_1 \equiv const. $ Для удобства, здесь и всюду далее положим $ \psi_2^0 = \psi_2(0). $
	
\subsection{Случай, когда $ u_1 \geqslant 0, \; u_2 \in [0, k], k > 0 $}
	\subsubsection{Ситуация, когда $ \psi_0 \neq 0 $}
	\label{fisrtCase1}
	Для того чтобы избавиться от свободы выбора $ \psi^{\star}(\cdot), $ рассмотрим ``нормальный'' случай, т.е $ \psi_0 < 0. $ Без ограничения общности положим $ \psi_0 = -\dfrac{1}{2}. $ Из условия максимума определим оптимальное управление $ u^{\star}(t) = (u_1^{\star}(t), u_2^{\star}(t)). $
	\begin{equation}
		\label{u1_Star}
		u_1^{\star} = 
		\begin{cases}
			\psi_2, & \psi_2 \geqslant 0, \\
			0, & \psi_2 < 0, \\   		
		\end{cases}
   	\end{equation}
   	\begin{equation}
   		\label{u2_Star}
   		u_2^{\star} = 
		\begin{cases}
			0, & \psi_2 x_2 > 0, \\
			[0, k], & \psi_2 x_2 = 0, \\ 
			k, & \psi_2 x_2 < 0.
   		\end{cases}
	\end{equation} 
	
	\begin{enumerate}
		\item Предположим, что $ \psi_2^0 > 0. $ Из  непрерывности функции $ \psi_2(t) $ следует, что  $ \exists \; \delta > 0: \psi_2(t) > 0, \; \forall t \in (0, \delta). $ Ввиду системы \eqref{u1_Star}: $ u_1^{\star}(t) = \psi_2(t) \; \forall t \in (0, \delta). $ Из формулы \eqref{x2(t)} следует, что $ x_2(t) > 0, \; \forall t \in (0, \delta). $ Значит из системы \eqref{u1_Star} следует, что $ u_2^{\star}(t) = 0, \; \forall t \in (0, \delta). $ \\
		Теперь, зная $ u_1^{\star}(t), u_2^{\star}(t) $, можно записать новую систему, состощую из вторых уравнений систем \eqref{odeSystem} и \eqref{conjSystem}
		\begin{equation}
		\label{firstSystem}
			\begin{cases}
				\dot{x_2} = \psi_2 - x_2, \\
				\dot{\psi_2} = -\psi_1 + \psi_2.   		
			\end{cases}
   		\end{equation}
   		Решим эту систему. Начнём со второго уравнения
   		$$
   			\dot{\psi_2}(t) = -\psi_1 + \psi_2(t).
   		$$
   		Решая данное уравнение методом вариации постоянной, получим
   		\begin{gather*}
			\psi_2(t) = \psi_1 + \tilde{C} e^t.   			
   		\end{gather*}
   		Константу $ \tilde{C} $ найдём из краевого условия: $ \psi_2(0) = \psi_2^0, $ в результате получим
   		\begin{gather}
		\label{psi2}
			\psi_2(t) = \psi_1 + (\psi_2^0 - \psi_1) e^t.
   		\end{gather}
   		
%%%%%%% Search x2(), x1(). %%%%%

		Для решения уравнения 
			$$
				\dot{x_2}(t) = \psi_2(t) - x_2(t),
			$$   			
тоже воспользуемся методом вариации постоянной. Решая, получим
   			\begin{gather*}
				x_2(t) = \psi_1 + \frac{1}{2}(\psi_2^0 - \psi_1) e^{t} + \tilde{C} e^{-t}.   			
   			\end{gather*}
   			Для нахождения константы $ \tilde{C}, $ воспользуемся краевым условием для системы \eqref{odeSystem}: $ x_2(0) = 0. $ В результате получим
   			\begin{gather}
   			\label{x2}
				x_2(t) = \psi_1(1 - e^{-t}) + (\psi_2^0 - \psi_1) \sh(t).
   			\end{gather}
   			И так, теперь, подставив $ x_2(t) $ в первое уравнение системы \eqref{odeSystem}, найдём $ x_1(t). $
   			\begin{gather*}
				\dot{x_1}(t) = \psi_1(1 - e^{-t}) + (\psi_2^0 - \psi_1) \sh(t) \\
				x_1(t) = \psi_1(t + e^{-t}) + (\psi_2^0 - \psi_1) \ch(t) + \tilde{C}.
   			\end{gather*}
			Подставляем краевое условие для системы \eqref{odeSystem}: $ x_1(0) = \alpha $ и получаем уравнение для $ x_1(t) $
			\begin{gather}
			\label{x1}
				x_1(t) = \psi_1(t + e^{-t}) + (\psi_2^0 - \psi_1) \ch(t) + \alpha - \psi_2^0.
   			\end{gather}
   			Для нахождения $ \psi_2, $ воспользуемся краевым условием из системы \eqref{odeSystem}: $ x_1(T) = L $
   			\begin{gather*}
				\psi_1(T + e^{-T}) + (\psi_2^0 - \psi_1) \ch(T) - \psi_2^0 = L \\
				\psi_2^0 = \frac{L - \alpha + \psi_1( \sh(T) - T)}{\ch(T) - 1}.
   			\end{gather*}
   			Воспользовавшись краевым условием для системы \eqref{odeSystem}: $ \vert x_2(T) \vert \leqslant \varepsilon, $ найдём ограничение на $ \psi_1. $ 
   			$$
   				-(L - \alpha) \sh(T) \leqslant \psi_1 \cdot f(T) \leqslant \varepsilon(\ch(T) - 1) - (L - \alpha) \sh(T), 
   			$$
   			где $ f(T) = (1 - e^{-T})(\ch(T) - 1) + \sh(T)(1 - T - e^{-T}). $
   			\begin{itemize}
   				\item Если $ f(T) > 0, $ тогда 
   				\begin{gather}
   				\label{psi1_bound_1}
   					\dfrac{-(L - \alpha) \sh(T)}{f(T)} \leqslant \psi_1 \leqslant \dfrac{\varepsilon(\ch(T) - 1) - (L - \alpha) \sh(T)}{f(T)};
   				\end{gather}
   				\item Если $ f(T) < 0, $ тогда 
   				\begin{gather}
   				\label{psi1_bound_2}
   					\dfrac{\varepsilon(\ch(T) - 1) - (L - \alpha) \sh(T)}{f(T)} \leqslant \psi_1 \leqslant \dfrac{-(L - \alpha) \sh(T)}{f(T)}.
   				\end{gather}
   			\end{itemize}
   			Далее будем исследоватть поведение функции $ \psi_2(t), $ при различных $ \psi_1, \psi_2^0. $
			\begin{enumerate}
			\item Предположим, что $ \psi_2^0 - \psi_1 > 0. $ Найдём производную функции $ \psi_2(t) $
			$$
				\psi_2'(t) = (\psi_2^0 - \psi_1) e^t.
			$$
			Сразу становится видно, что $ \psi_2'(t) > 0, \; \forall t \in [0, T]. $ Это значит, что функция $ \psi_2(t) $ строго возрастает при $ t \in [0, T]. $ Следовательно $ \psi_2(t) > 0, \; \forall t \in [0, T]. $ В силу \eqref{u1_Star}: $ u_1^{\star}(t) = \psi_2(t), \; \forall t \in [0, T]. $ \\
			Из формулы \eqref{x2(t)} видно: $ x_2(t) > 0, \; \forall t \in (0, T]. $ Ну а раз так, то в силу \eqref{u2_Star}: $ u_2^{\star}(t) = 0, \; \forall t \in (0, T]. $
   			\item Пусть теперь $ \psi_1 - \psi_2^0 = 0. $ Тогда получается, что $ \psi_2(t) > 0, \; \forall t \in [0, T]. $ В силу \eqref{u1_Star}: $ u_1^{\star}(t) = \psi_2(t), \; \forall t \in [0, T]. $ \\
			Из формулы \eqref{x2(t)} видно: $ x_2(t) > 0, \; \forall t \in (0, T]. $ Ну а раз так, то в силу \eqref{u2_Star}: $ u_2^{\star}(t) = 0, \; \forall t \in (0, T]. $
			
   			Перепишем уравнения \eqref{psi2}, \eqref{x2}, \eqref{x1}
   			\begin{gather*}
   				\begin{cases}
   					\psi_2(t) = \psi_2^0, \\
   					x_2(t) = \psi_2^0(1 - e^{-t}), \\
   					x_1(t) = \psi_2^0(t + e^{-t} - 1) + \alpha,
   				\end{cases}  			
   			\end{gather*}
   			с краевыми условиями
   			\begin{gather*}
   				\begin{cases}
   					\vert x_2(T) \vert = \psi_2^0(1 - e^{-T}) \leqslant \varepsilon, \\
   					x_1(T) = \psi_2^0(T + e^{-T} - 1) + \alpha = L.
   				\end{cases}
   			\end{gather*}
   			Из второго краевого условия получаем
   			$$
   				\psi_2^0 = \dfrac{L - \alpha}{T + e^{-T} - 1},
   			$$
   			а из первого
   			$$
   				0 \leqslant \dfrac{(L - \alpha)(1 - e^{-T})}{T + e^{-T} - 1} \leqslant \varepsilon.
   			$$
   			
   			\item \label{firstCase_c}
   			Пусть $ \psi_2^0 - \psi_1 < 0. $ Вычислим производную функции $ \psi_2(t) $ 
   			$$
				\psi_2'(t) = (\psi_2^0 - \psi_1) e^t.
			$$
			Сразу становится видно, что $ \psi_2'(t) < 0, \; \forall t \in [0, T]. $ Это значит, что функция $ \psi_2(t) $ строго убывает на всем промежутке времени $ [0, T]. $ А раз так, то $ \exists \; t_{\text{пер}} = \ln \left( -\dfrac{\psi_1}{\psi_2^0 - \psi_1} \right) > 0 : \psi_2(t_{\text{пер}}) = 0. $ \\
			Теперь рассмотрим время, до момента переключения и после.
			\begin{itemize}
				\item Если $ t \in [0, t_{\text{пер}}), $ то $ \psi_2(t) > 0. $ Из формулы \eqref{x2(t)} следует, что $ x_2(t) > 0. $ Теперь мы можем выписать управление из систем \eqref{u1_Star}, \eqref{u2_Star}: $ u_1^{\star}(t) = \psi_2(t), u_2^{\star}(t) = 0. $
				\item Если $ t \in [t_{\text{пер}}, T], $ то $ \psi_2(t) < 0. $ Из формулы \eqref{x2} следует, что $ x_2(t) > 0. $ Теперь можно выписать систему
				\begin{gather*}
   					\begin{cases}
   						\dot{\psi}_2 = -\psi_1 + \psi_2(1 + k), \\
   						\dot{x}_2 = - x_2(1 + k), \\
   						\dot{x}_1 = x_2, 
					\end{cases}  			
   				\end{gather*}
   		с начальными условиями
   				\begin{gather*}
   					\begin{cases}
   						x_2(t_{\text{пер}}) = \psi_1(1 - e^{-t_{\text{пер}}}) + (\psi_2^0 - \psi_1) \sh(t_{\text{пер}}), \\
   						x_1(t_{\text{пер}}) = \psi_1(t_{\text{пер}} + e^{-t_{\text{пер}}}) + (\psi_2^0 - \psi_1) \ch(t_{\text{пер}}) + \alpha - \psi_2^0, \\
   						\psi_2(t_{\text{пер}}) =0.   							\end{cases}  			
   				\end{gather*}
   				Решая эту систему, получаем
   				$$
   					\psi_2(t) = \dfrac{\psi_1}{1 + k} \left(1 - e^{(1 + k)(t - t_{\text{пер}})} \right),
   				$$
   				$$
   					x_2(t) =  \left( \psi_1(1 - e^{-t_{\text{пер}}}) + (\psi_2^0 - \psi_1) \sh(t_{\text{пер}})\right) e^{-(1 + k)(t - t_{\text{пер}})}.
   				$$
   				\begin{gather*}
   					x_1(t) =  \dfrac{1}{1 + k}[\psi_1(1 - e^{-t_{\text{пер}}}) + (\psi_2^0 - \psi_1) \sh(t_{\text{пер}})](1 - e^{-(1 + k)(t - t_{\text{пер}})}) + \\ + \psi_1(t_{\text{пер}} + e^{-t_{\text{пер}}}) + (\psi_2^0 - \psi_1) \ch(t_{\text{пер}}) + \alpha - \psi_2^0.
   				\end{gather*}
   				Из формулы \eqref{psi2} можно выразить $ \psi_1 $ через $ \psi_2 $
   				$$
   					\psi_1 = \dfrac{\psi_2^0 e^{t_{\text{пер}}}}{e^{t_{\text{пер}}} - 1}.
   				$$
   				Подставляя $ \psi_1 $ в $ x_1(t), $ окончательно получим
   				\begin{gather*}
   					x_1(t) =  \dfrac{\psi_2^0}{e^{t_{\text{пер}}} - 1} \left(\dfrac{1}{1 + k}[e^{t_{\text{пер}}} - 1 - \sh(t_{\text{пер}})](1 - e^{-(1 + k)(t - t_{\text{пер}})}) - \ch(t_{\text{пер}}) + \right. \\ \left. + e^{t_{\text{пер}}}(t_{\text{пер}} - 1) + 2 \right) + \alpha.
   				\end{gather*}
   				Введём обозначение
   				\begin{gather*}
   					F(t_{\text{пер}}) = \dfrac{1}{1 + k}[e^{t_{\text{пер}}} - 1 - \sh(t_{\text{пер}})](1 - e^{-(1 + k)(T - t_{\text{пер}})}) - \ch(t_{\text{пер}}) + \\ + e^{t_{\text{пер}}}(t_{\text{пер}} - 1) + 2.
   				\end{gather*}
   				Тогда из условия: $ x_1(T) = L $ следует
   				$$
   					\psi_2^0 = \dfrac{(L - \alpha)(e^{t_{\text{пер}}} - 1)}{F(t_{\text{пер}})}.
   				$$
			\end{itemize}
   		\end{enumerate}
   	\item Пусть $ \psi_2^0 = 0. $ Тогда управление определено неоднозначно. Но можно заметить, что 
   	$$ 
   		\dot{\psi_2}(0) =  - \psi_1.
   	$$
   	\begin{itemize}
   		\item Если $ \psi_1 < 0. $ Тогда $ \exists \; \delta > 0 : \psi_2(t) > 0, \; \forall t \in (0, \delta). $ В силу формулы \eqref{x2(t)}: $ x_2(t) > 0, \; \forall t \in (0, \delta). $ Учитывая системы \eqref{u1_Star}, \eqref{u2_Star}, можно выписать оптимальное управление $ u_1^{\star}(t) = \psi_2(t), u_2^{\star}(t) = 0, \; \forall t \in (0, \delta). $ \\
		Тогда будет справедлива система \eqref{firstSystem}.
   		Как было показано выше, данная система при $ \psi_2^0 = 0 $ имеет следующее решение
   		\begin{gather*}
   				\begin{cases}
   					\psi_2(t) = \psi_1(1 - e^t), \\
   					x_2(t) = \psi_1(1 - \ch(t)), \\
   					x_1(t) = \psi_1(t - \sh(t)) + \alpha.
   				\end{cases}  			
   		\end{gather*}
   		Можно заметить, что $ \psi_2(t) > 0, \; \forall t \in (0, T]. $
   		Воспользовавшись краевым условием: $ x_1(T) = L, $ получаем
   		$$
   			\psi_1 = \frac{L - \alpha}{T - \sh(T)}.
   		$$
   		А из краевого условия: $ \vert x_2(T) \vert \leqslant \varepsilon, $ получаем ограничение на параметры $ L, T $
   		$$
   			0 \leqslant \frac{(L - \alpha)(1 - \ch(T))}{T - \sh(T)} \leqslant \varepsilon.
   		$$
   		\item Если $ \psi_1 > 0. $ Тогда $ \exists \; \delta > 0 : \psi_2(t) < 0, \; \forall t \in (0, \delta). $ В силу формулы \eqref{x2(t)}: $ x_2(t) = 0, \; \forall t \in (0, \delta). $ Как можно заметить, $ \dot{\psi_2}( t_{\text{нач}}) < 0. $ В этом случае мы сможем добиться того, чтобы $ \psi_2(t) > 0. $ Следовательно $ x_2(t) = 0, \; \forall t \in [0, T]. $
		\item Если $ \psi_1 = 0. $ В этом случае получаем противоречие с условием нетривиальности вектора сопряжённых переменных. 
   	\end{itemize}
   	\item Пусть $ \psi_2^0 < 0. $ \\
   	\label{u1PosPsi_2^0<0}
   	В силу непрерывности $ \psi_2(t) < 0, \; \forall t  \in (0, \delta). $ И системы \eqref{u1_Star} cледует, что $ u_1^{\star}(t) =  0, \; \forall t \in (0, \delta). $ Тогда по формуле \eqref{x2(t)} получается, что $ x_2(t) =  0, \; \forall t \in (0, \delta). $ Предположим, что $ \exists \; t_{\text{нач}} > 0 : \psi_2(t_{\text{нач}}) = 0. $ Если мы вычислим производную функции $ \psi_2(t) $ в точке $ t_{\text{нач}}, $ то увидим, что
	\begin{enumerate}
		\item если $ \psi_1 > 0, $ тогда $ \dot{\psi_2}(t_{\text{нач}}) < 0. $ Этот случай не интересен, так как $ \dot{\psi_2}(t) $ в момент времени $ t = t_{\text{нач}} $ всегда отрицательна. Это значит, что $ x_2(t) = 0, \; \forall t \in [0, T].  $
		\item если $ \psi_1 = 0, $ тогда возникает противоречие с нетривиальностью вектора сопряжённых переменных.
		\item если $ \psi_1 < 0, $ тогда $ \dot{\psi_2}(t_{\text{нач}}) > 0. $ Это значит, что $ \exists \; \delta > 0 : \psi_2(t) > 0, \; \forall t \in (t_{\text{нач}}, \delta). $ В силу системы \eqref{u1_Star}, \eqref{u2_Star}: $ u_1^{\star} = \psi_2(t), u_2^{\star} = k. \; \forall t \in (t_{\text{нач}}, \delta). $ Теперь мы можем выписать новую систему
		\begin{gather*}
   				\begin{cases}
   					\dot{\psi_2} = - \psi_1 + \psi_2(1 + k), \\
   					\dot{x}_2 = \psi_2 - x_2(1 + k). \\
   				\end{cases}
   		\end{gather*}   	
   		Решим первое уравнение этой системы с краевым условием: $ \psi_2(t_{\text{нач}}) = 0 $
   		\begin{gather*}
   			\psi_2(t) = \dfrac{\psi_1}{1 + k} \left( 1 - e^{(1 + k) (t - t_{\text{нач}})} \right).
   		\end{gather*}
   		Как можно видеть, $ \psi_2(t) > 0, \; \forall t \in (t_{\text{нач}}, T]. $ Ввиду формулы \eqref{x2(t)}: $ x_2(t) > 0, \; \forall t \in (t_{\text{нач}}, T]. $
   		Решим второе уравнение этой системы с краевым условием: $ x_2(t_{\text{нач}}) = 0 $
   		\begin{gather*}
   			x_2(t) = \dfrac{\psi_1}{(1 + k)^2} [ 1 -  \ch((1 + k)(t - t_{\text{нач}}))].
   		\end{gather*}
   		Зная $ x_2(t) $ и краевое условие: $ x_1(t_{\text{нач}}) = \alpha, $ можно найти $ x_1(t) $
   		\begin{gather*}
   			x_1(t) = \dfrac{\psi_1}{(1 + k)^2} \left[ t - t_{\text{нач}} -  \dfrac{1}{1 + k}\sh((1 + k)(t - t_{\text{нач}})) \right] + \alpha.
   		\end{gather*}
   		Из краевого условия: $ x_1(T) = L, $ можно найти $ \psi_1 $
   		\begin{gather*}
   			\psi_1 = \dfrac{(L - a) (1 + k)^3}{(T - t_{\text{нач}})(1 + k) - \sh((1 + k)(T - t_{\text{нач}}))}.
   		\end{gather*}
   		Вспоминая краевое условие: $ x_2(T) \leqslant \varepsilon, $ получим ограничение на $ t_{\text{нач}}. $
   		\end{enumerate}	 		
	\end{enumerate}
	
	\subsubsection{Случай, когда $ \psi_0 = 0 $ }
	Теперь расмотрим случай, когда $ \psi_0 = 0. $
	Выпешем из условия максимума оптимальное управление $ u^{\star}(t) = (u_1^{\star}(t), u_2^{\star}(t)). $
	\begin{equation*}
		u_1^{\star} = 
		\begin{cases}
			0, & \psi_2 < 0, \\
			[0, +\infty), & \psi_2 = 0, \\ 
			+\infty, & \psi_2 > 0,
   		\end{cases}
   	\end{equation*}
   	\begin{equation*}
   		u_2^{\star} = 
		\begin{cases}
			0, & \psi_2 x_2 > 0, \\
			[0, k], & \psi_2 x_2 = 0, \\ 
			k, & \psi_2 x_2 < 0.
   		\end{cases}
	\end{equation*}	

	\begin{enumerate}
		\item Случай $ \psi_2(t) > 0 $ мы не рассматриваем, ввиду того, что $ u_1(t) $ становится неограниченным.
		\item Пусть $ \psi_2^0 < 0. $ Ввиду непрерывности функции $ \psi_2(t), \exists \; \delta > 0 : \psi_2(t) < 0, \; \forall t \in (0, \delta). $ Пердположим, что $ \exists \; t_{\text{нач}} > 0 : \psi_2(t_{\text{нач}}) = 0. $ И так, рассмотрим случай до мемента переключения и после.
		\begin{enumerate}
			\item Если $ t \in [0, t_{\text{нач}}), $ то $ \psi_2(t) < 0. $ По формуле \eqref{x2(t)}: $ x_2(t) = 0. $
			\item Если $ t = t_{\text{нач}}, $ то посмотрим на производную $ \dot{\psi_2}(t) = - \psi_1. $ Видно, что
			\begin{itemize}						
			\item если $ \psi_1 > 0, $ то $ \dot{\psi_2}(t) < 0, $ а значит $ \exists \; \delta > 0 : \psi_2(t) < 0 \Rightarrow x_2(t) = 0. $
			\item если $ \psi_1 < 0, $ то $ \dot{\psi_2}(t) > 0, $ а значит $ \exists \; \delta > 0 : \psi_2(t) > 0. $ Этот случай нам не интересен, поскольку управление не ограничено.
			\item если $ \psi_1 = 0, $ то возникает противоречие с нетривиальностью вектора сопряжённых переменных.
			\end{itemize}
		\end{enumerate}
		И так, в данном случае не существует оптимального оправления.
		\item Пусть $ \psi_2^0 = 0. $ Значит $ \dot{\psi_2}(0) = - \psi_1. $
		\begin{itemize}
			\item Если $ \psi_1 = 0, $ то возникает противоречие с нетривиальностью вектора сопряжённых переменных.
			\item Рассмотрим случай, когда $ \psi_1 > 0. $ \\ 
			В данном случае $ \dot{\psi_2}(0) = -\psi_1 \Rightarrow \dot{\psi_2}(0) < 0. $ Следовательно $ \exists \; \delta > 0 : \psi_2(t) < 0, \forall t \in (0, \delta). $ Мы получаем, что $ x_2(t) = 0, \;  \forall t \in (0, \delta). $ 	
			\item Рассмотрим случай, когда $ \psi_1 < 0. $ \\
			В данном случае $ \dot{\psi_2}(0) = -\psi_1 \Rightarrow \dot{\psi_2}(0) > 0. $ Следовательно $ \exists \; \delta > 0 : \psi_2(t) > 0, \forall t \in (0, \delta). $ Мы не рассматриваем данный случай, так как управление не ограничено.
		\end{itemize}
		Получается, что в данном случае тоже не существует оптимального управления.
	\end{enumerate}	
\subsection{Случай, когда $ u_1 \in \mathbb{R}, u_2 \in [0, k], \; k > 0 $}	
	\subsubsection{Случай, когда $ \psi_0 \neq 0 $}
	Как и ранее рассмотрим $ \psi_0 = - \dfrac{1}{2}. $ Из условия максимума определим управление, $ u^{\star}(t) = (u_1^{\star}(t), u_2^{\star}(t)), $ подозреваемое на оптимальность.
	\begin{gather}
	\label{secondCase_u1u2}
		u_1^{\star} = \psi_2(t), \; \;
		u_2^{\star} = 
   		\begin{cases}
   			0, & \psi_2 x_2 > 0, \\
			[0, k], & \psi_2 x_2 = 0, \\ 
			k, & \psi_2 x_2 < 0.
   		\end{cases}  			
   	\end{gather}
   	Можно заметить, вектор управления будет выглядеть следующим образом $ u^{\star}(t) = (\psi_2(t), 0). $ Это вытекает из формулы \eqref{x2(t)}, а именно то, что $ \psi_2(t) $ и $ x_2(t) $ всегда будут отдного знака, за исключением одного случая, когда $ \psi_2(t) = 0. $ В этом случае вектор управления определяется неоднозначно.
   	\begin{enumerate}
   		\item Пердположим, что $ \psi_2^0 > 0. $ Данный случай полностью соответсвует случаю, когда $ \psi_2^0 > 0, $ в \eqref{fisrtCase1}, за исключением пункта \eqref{firstCase_c} при $ t \in (t_{\text{пер}}, T]. $ \\
   		В этом случае $ \psi_2(t) < 0. $ Из формулы \eqref{x2(t)} следует, что $ x_2(t) < 0. $ Поэтому, в силу системы \eqref{secondCase_u1u2}, можно записать следующую систему
   		\begin{gather*}
   			\begin{cases}
   				\dot{x_2} = \psi_2  - x_2, \\
   				\dot{\psi_2} = -\psi_1 + \psi_2,
   			\end{cases}  			
   		\end{gather*}
   	с краевыми условиями
   	\begin{gather*}
   			\begin{cases}
   				x_2(t_{\text{пер}}) = \psi_1(1 - e^{-t_{\text{пер}}}) + (\psi_2^0 - \psi_1) \sh(t_{\text{пер}}), \\
   				\psi_2(t_{\text{пер}}) = 0.
   			\end{cases}  			
   		\end{gather*}
   	Решим эту систему. Начнём со второго уравнения
   		$$
   			\dot{\psi_2}(t) = -\psi_1 + \psi_2(t).
   		$$
   		Решая данное уравнение методом вариации постоянной, получим
   		\begin{gather*}
			\psi_2(t) = \psi_1 + \tilde{C} e^t.   			
   		\end{gather*}
   		Константу $ \tilde{C} $ найдём из краевого условия: $ \psi_2(t_{\text{пер}}) = 0, $ в результате получим
   		\begin{gather*}
			\psi_2(t) = \psi_1(e^{t - t_{\text{пер}}} - 1).
   		\end{gather*}
   		
%%%%%%% Search x2(), x1(). %%%%%

		Тогда, для решения уравнения 
			$$
				\dot{x_2}(t) = \psi_2(t) - x_2(t),
			$$   			
воспользуемся методом вариации постоянной. Решая, получим
   			\begin{gather*}
				x_2(t) = \psi_1 + \frac{1}{2}(\psi_2^0 - \psi_1) e^{t} + \tilde{C} e^{-t}.   			
   			\end{gather*}
   			Для нахождения константы $ \tilde{C}, $ воспользуемся краевым условием: $ x_2(t_{\text{пер}}) = \psi_1(1 - e^{-t_{\text{пер}}}) + ( \psi_2^0 - \psi_1) \sh(t_{\text{пер}}). $ В результате получим
   			\begin{gather*}
				x_2(t) = \dfrac{\psi_1}{2}e^{t - t_{\text{пер}}} - \psi_1 + \left(\dfrac{3}{2} \psi_1 - \psi_1 e^{-t_{\text{пер}}} + ( \psi_2^0 - \psi_1) \sh(t_{\text{пер}}) \right) e^{t_{\text{пер}} - t}.
   			\end{gather*}
   			И так, используя краевое условие: $ x_1(t_{\text{пер}}) = \psi_1(t_{\text{пер}} - e^{-t_{\text{пер}}}) + (\psi_2^0 - \psi_1) \ch(t_{\text{пер}}) + \alpha - \psi_2^0$, найдём $ x_1(t). $
   			\begin{gather*}
				x_1(t) = \dfrac{\psi_1}{2}e^{t - t_{\text{пер}}} - \psi_1 t - \left( \dfrac{3}{2}\psi_1 - \psi_1 e^{-t_{\text{пер}}} + (\psi_2^0 - \psi_1) \sh(t_{\text{пер}}) \right) e^{t_{\text{пер}} - t} + \\ + 2\psi_1(t_{\text{пер}} - e^{-t_{\text{пер}}}) + (\psi_2^0 - \psi_1)\left(e^{-t_{\text{пер}}} - 1 \right) + \alpha.
   			\end{gather*}
   			Для нахождения $ \psi_1, $ воспользуемся краевым условием из системы \eqref{odeSystem}: $ x_1(T) = L $
   			\begin{gather*}
				\psi_1 = \dfrac{L + \psi_2^0(e^{-t_{\text{пер}}} - 1 - \sh(t_{\text{пер}})e^{t_{\text{пер}} - T}) + \alpha}{\sh(T - t_{\text{пер}}) - T + e^{-T} + (\sh(t_{\text{пер}}) - 1)e^{t_{\text{пер}} - T} + 2 t_{\text{пер}} - 3e^{-t_{\text{пер}}} + 1}.
   			\end{gather*}
   			Воспользовавшись краевым условием для системы \eqref{odeSystem}: $ \vert x_2(T) \vert \leqslant \varepsilon, $ найдём ограничение на $ t_{\text{пер}}. $
   		\item Пусть $ \psi_2^0 < 0. $
   		Из непрерывности $ \psi_2(t) $ следует, что $ \psi_2(t) < 0, \; \forall t \in (0, \delta). $ Из формулы \eqref{x2(t)} следует, что $ x_2(t) < 0, \; \forall t \in (0, \delta). $ Получается следующий вектор управления $ u^{\star}(t) = (\psi_2(t), 0) $ и новая система
   		\begin{gather*}
   			\begin{cases}
   				\dot{x_2} = \psi_2 - x_2, \\
   				\dot{\psi_2} = -\psi_1 + \psi_2.
   			\end{cases}  			
   		\end{gather*}
   		Для данной системы справедливы формулы \eqref{psi2}, \eqref{x2}, \eqref{x1}.
   		\begin{enumerate}
			\item Предположим, что $ \psi_2^0 - \psi_1 > 0. $ Найдём производную функции $ \psi_2(t) $
			$$
				\psi_2'(t) = (\psi_2^0 - \psi_1) e^t.
			$$
			Сразу становится видно, что $ \psi_2'(t) > 0, \; \forall t \in [0, T]. $ Это значит, что функция $ \psi_2(t) $ строго возрастает при $ t \in [0, T]. $ Следовательно $ \exists \; t_{\text{пер}} > 0 : \psi_2(t_{\text{пер}}) = 0. $ Ввиду системы \eqref{secondCase_u1u2}, вектор управления, после момента переключения, останется тем же, т.е. $ u^{\star}(t) = (\psi_2(t), 0). $
   			\item Пусть теперь $ \psi_1 - \psi_2^0 = 0. $ Тогда получается, что $ \psi_2(t) < 0, \; \forall t \in [0, T]. $ В силу \eqref{secondCase_u1u2}: $ u_1^{\star}(t) = \psi_2(t), \; \forall t \in [0, T]. $ \\
			Из формулы \eqref{x2(t)} видно: $ x_2(t) < 0, \; \forall t \in [0, T]. $ Ну а раз так, то в силу \eqref{secondCase_u1u2}: $ u_2^{\star}(t) = 0, \; \forall t \in [0, T]. $ Следовательно никаких переключений не будет на $ [0, T]. $
   			\item Пусть $ \psi_2^0 - \psi_1 < 0. $ Вычислим производную функции $ \psi_2(t) $ 
   			$$
				\psi_2'(t) = (\psi_2^0 - \psi_1) e^t.
			$$
			Сразу становится видно, что $ \psi_2'(t) < 0, \; \forall t \in [0, T]. $ Это значит, что функция $ \psi_2(t) $ строго убывает на всем промежутке времени $ [0, T]. $ Следовательно никаких переключений не будет на $ [0, T]. $
   		\end{enumerate}
   		 \item Пусть $ \psi_2^0 = 0. $ Посмотрим производную функции $ \psi_2(t) $ в нуле
   		$$
			\psi_2'(0) = - \psi_1.
		$$
		Возножны три случая:
		\begin{enumerate}
			\item если $ \psi_1 = 0, $ тогда возникает противоречие с нетривиальностью вектора сопряжённых переменных.
			\item если $ \psi_1 > 0, $ тогда $ \exists \; \delta > 0 : \psi_2(t) < 0, \; \forall t \in (0, \delta). $ А раз так, то по формуле \eqref{secondCase_u1u2}: $ x_2(t) < 0, \; \forall t \in (0, \delta). $ Выпишем вектор управления $ u^{\star}(t) = (\psi_2(t), 0) $ и систему, отвечающую этому управлению 
		\begin{gather*}
   			\begin{cases}
   				\dot{x_2} = \psi_2 - x_2, \\
   				\dot{\psi_2} = -\psi_1 + \psi_2.
   			\end{cases}  			
   		\end{gather*}
   		Для данной системы справедливы формулы \eqref{psi2}, \eqref{x2}, \eqref{x1}.
			Найдём производную функции $ \psi_2(t) $
			$$
				\psi_2'(t) = - \psi_1 e^t.
			$$
			Сразу становится видно, что $ \psi_2'(t) < 0, \; \forall t \in [0, T]. $ Это значит, что функция $ \psi_2(t) $ строго убывает при $ t \in [0, T]. $ Следовательно никаких переключений не будет на $ [0, T]. $
			\item если $ \psi_1 < 0, $ тогда $ \exists \; \delta > 0 : \psi_2(t) > 0, \; \forall t \in (0, \delta). $ А раз так, то по формуле \eqref{secondCase_u1u2}: $ x_2(t) > 0, \; \forall t \in (0, \delta). $ Выпишем вектор управления $ u^{\star}(t) = (\psi_2(t), 0) $ и систему, отвечающую этому управлению 
		\begin{gather*}
   			\begin{cases}
   				\dot{x_2} = \psi_2 - x_2, \\
   				\dot{\psi_2} = -\psi_1 + \psi_2.
   			\end{cases}  			
   		\end{gather*}
   		Для данной системы справедливы формулы \eqref{psi2}, \eqref{x2}, \eqref{x1}.
			Найдём производную функции $ \psi_2(t) $
			$$
				\psi_2'(t) = - \psi_1 e^t.
			$$
			Сразу становится видно, что $ \psi_2'(t) > 0, \; \forall t \in [0, T]. $ Это значит, что функция $ \psi_2(t) $ строго возрастает при $ t \in [0, T]. $ Следовательно никаких переключений не будет на $ [0, T]. $
		\end{enumerate}
   	\end{enumerate}
   	\subsubsection{Случай, когда $ \psi_0 = 0 $}
	Выпешем из условия максимума оптимальное управление $ u^{\star}(t) = (u_1^{\star}(t), u_2^{\star}(t)). $
	\begin{equation*}
		u_1^{\star} = 
		\begin{cases}
			-\infty, & \psi_2 < 0, \\
			\mathbb{R}, & \psi_2 = 0, \\ 
			+\infty, & \psi_2 > 0,
   		\end{cases}
   	\end{equation*}
   	\begin{equation*}
   		u_2^{\star} = 
		\begin{cases}
			0, & \psi_2 x_2 > 0, \\
			[0, k], & \psi_2 x_2 = 0, \\ 
			k, & \psi_2 x_2 < 0.
   		\end{cases}
	\end{equation*}
	Нас интересуют только конечные управления, следовательно $ u_1^{\star} = \mathbb{R} \Leftrightarrow \psi_2  \equiv 0. $
	В результате управление не определено, и следовательно наша система не является управляемой.
\section{Выводы}
	И так, чтобы вычислить оптимальное управление, нужно перебрать все управления, подозреваемые на оптимальность и посмотреть, какое из них минимизурет функционал $ J = \int \limits_0^T u_1^2(t)dt \rightarrow \min \limits_{u(\cdot)}. $ 
\section{Особенности реализации программы}
	В некоторых случаях нам нужно оценить время переключния. Для этого нужно решить трансцендентное уравнение. Для решения используется функция \texttt{fzero(fun, [x0, x1])}. \\
	В формулах \eqref{psi1_bound_1}, \eqref{psi1_bound_2} мы перебераем по $ \psi_1 $ из соответствующего промежутка и подставляем в формулы для $ x_2(t), x_1(t). $ 
\end{document}
